
\section{Background}

\subsection{Softmax Attention}

For $Q, K, V \in \mathbb{R}^{L \times d_h}$:

\[
A = \text{softmax}\!\left(\frac{Q K^\top}{\sqrt{d_h}}\right) \in \mathbb{R}^{L \times L}
\]

\[
\text{Attn}(Q, K, V) = A V \in \mathbb{R}^{L \times d_h}
\]

\textbf{Complexity.} $A$ requires $\mathcal{O}(L^2)$ memory; the two matmuls 
require $\mathcal{O}(L^2 d_h)$ time.

\subsection{Nystr m Approximation}

For a symmetric PSD kernel matrix $K \in \mathbb{R}^{L \times L}$, the Nystr m 
method selects $m$ landmark rows/columns and approximates:

\[
K \approx C W^{-1} C^\top
\]

where:

\begin{itemize}
\item $C \in \mathbb{R}^{L \times m}$: query-landmark kernel
\item $W \in \mathbb{R}^{m \times m}$: landmark-landmark kernel (PSD)

\end{itemize}
The approximation is exact when $K$ has rank $\le m$.

\textbf{Application to attention.} With $K_\text{attn} = \text{softmax}(QK^\top / \sqrt{d_h})$, 
the Nystr mform is:

\[
A \approx \tilde{A} = \begin{pmatrix} A_{Q, L_m} \\ A_{L_m, L_m} \\ A_{L_m, Q} \end{pmatrix}
\]

where $L_m$ indexes $m$ landmarks.

\subsection{Asymmetric Kernel}

When the kernel is \textbf{not symmetric} -- e.g., $K(q_i, k_j) \ne K(k_i, q_j)$ 
-- Nystr m approximation in the standard form does not apply.

\textbf{Observation.} The attention kernel $k(q, k) = \exp(q^\top k / \sqrt{d_h})$ 
\textbf{is} symmetric in $q, k$ for the \textit{score} function, but the matrices $Q$ and 
$K$ themselves contain different information. The \textbf{row and column spaces} 
of $A$ may differ.

Primal Attention handles this by allowing $Q_\text{land} \ne K_\text{land}$.

\subsection{Primal vs Dual Representation}

Given a kernel matrix $K$ with rank $r$:

\begin{itemize}
\item \textbf{Primal representation:} $K = \Phi \Phi^\top$ where $\Phi \in \mathbb{R}^{L \times r}$
\item \textbf{Dual representation:} $K = U \Sigma U^\top$ where $U \in \mathbb{R}^{L \times r}$ is orthogonal

\end{itemize}
Primal is cheaper when $r \ll L$ and $K$ is low-rank. Primal Attention uses 
the primal form with asymmetric factors.

\subsection{Moore--Penrose Pseudoinverse}

For a matrix $A \in \mathbb{R}^{m \times n}$, the Moore--Penrose pseudoinverse 
$A^+ \in \mathbb{R}^{n \times m}$ satisfies:

\[
A A^+ A = A, \qquad A^+ A A^+ = A^+, \qquad (A A^+)^\top = A A^+, \qquad (A^+ A)^\top = A^+ A
\]

\textbf{Computation.} For $A = U \Sigma V^\top$, $A^+ = V \Sigma^+ U^\top$ where 
$\Sigma^+_{ii} = 1/\sigma_i$ if $\sigma_i > \epsilon$, else $0$.

\subsection{Schulz Iteration}

Schulz (1933) proposed the iteration:

\[
X_{t+1} = X_t (2I - A X_t)
\]

which converges to $A^{-1}$ \textbf{quadratically} if $\|I - A X_0\| < 1$.

\textbf{Convergence.} Let $E_t = I - A X_t$. Then:

\[
E_{t+1} = (I - A X_{t+1}) = (I - A X_t (2I - A X_t)) = E_t^2
\]

so $\|E_t\| \le \|E_0\|^{2^t}$, giving quadratic convergence.

\textbf{Initialization.} Common choices: $X_0 = A^\top / \|A\|_F^2$ or $X_0 = A^+$ 
(precomputed).

\subsection{LoRA Recovery}

LoRA (Hu et al., 2021) adds trainable low-rank adapters:

\[
W_\text{eff} = W_\text{frozen} + \frac{\alpha}{r_L} B A
\]

We use LoRA to recover quality after attention replacement.

\subsection{Notation Summary}

\begin{lstlisting}

---
\end{lstlisting}